% Ad Familiares 6.16 --- headnote
% ad-familiares-06-16, late March 44 BC, in Sicily

\textit{A.~Pompeius Bithynicus to Cicero, written from Sicily in late
March 44 BC, in the immediate aftermath of Caesar's assassination
(Perseus dateline: \textit{in Sicilia exeunte mense Martio 710
(44)}). This is the unusual letter in the cluster --- and one of
relatively few in the whole \textit{Familiares} --- in which Cicero
is the recipient, not the sender. The salutation
\textit{BITHYNICVS CICERONI S.} flags it immediately. Bithynicus,
serving as governor of Sicily, writes to ask Cicero to look after
his interests at Rome in his absence.}

\textit{The letter is short and turns on a graceful opening figure:
Bithynicus says he could appeal to the hereditary friendship between
his father (also A.~Pompeius Bithynicus, a long-time correspondent
of Cicero's) and Cicero --- but that, he says, is the move of men
who have themselves done nothing to deserve a friendship of their
own, and he means to stand on the one he has built with Cicero in
his own right. The request that follows is therefore made in his
own name. Cicero's reply is the next letter, \textit{Fam.}~6.17.}
